\section{The degreewise affine method}
\label{sec:method}

For a general $r\ge2$, define
\begin{equation}\label{eq:Fr}
 F_r(x)=1+x\cos s-\sqrt{1-x^2}\sin s.
\end{equation}
Then $H_d(r)=F_r(x_d)$.  On the interval $x_{11}\le x\le x_1$,
\begin{equation}\label{eq:convexity}
 F_r''(x)=\frac{\sin s}{(1-x^2)^{3/2}}>0
\end{equation}
for every $r\ge2$.  Thus $F_r$ lies below its endpoint chord.  Since $x_d$ is affine in $d$, the
barycentric weights are $(d-1)/10$ at $x_{11}$ and $(11-d)/10$ at $x_1$.
Define
\[
 M_1(r)=\frac{H_1(r)}{11},
 \qquad
 M_{11}(r)=H_{11}(r),
 \qquad
 M(r)=\max_{1\le d\le11}\frac{H_d(r)}{12-d}.
\]
The chord estimate gives
\begin{equation}\label{eq:endpointreduction}
 \boxed{M(r)=\max\{M_1(r),M_{11}(r)\}.}
\end{equation}
Indeed, the endpoint degrees attain the two ratios, while every intermediate
ratio is bounded by their maximum.

Define the exposed constituent-sphere set and its area by
\[
 X_i(r)=\partial(c_i+rB)\setminus
        \bigcup_{j\ne i}\operatorname{int}(c_j+rB),
 \qquad
 e_i(r)=\cH^2(X_i(r)).
\]
For $d_i\le11$,
\[
 e_i(r)\le2\pi r^2M(r)(12-d_i).
\]
For $d_i=12$, the twelve closed radius-$\pi/3$ caps cover $\Sph^2$ by the kissing
number theorem.  At $r=2$, the possible residue consists only of finitely many
sphere-intersection circles and has zero area; for $r>2$ the actual hidden caps are
strictly larger.  Thus $e_i(r)=0$ throughout $r\ge2$.  Every point of
$\partial U_r$ belongs to at least one $X_i(r)$.  Using \cref{eq:handshake},
\begin{equation}\label{eq:surfaceuppergeneral}
 \cH^2(\partial U_r)
 \le4\pi r^2M(r)D.
\end{equation}
Thus, whenever the finite global surface lower bound is available at $r$,
maximizing the contact coefficient is equivalent to minimizing $r^2M(r)$.

\section{Exact optimization of the radius}
\label{sec:optimizer}

Because $\rho_1=\beta$,
\[
 H_1(r)=1+\cos\alpha=1+\frac1r,
\]
so
\begin{equation}\label{eq:B1}
 B_1(r):=r^2M_1(r)=\frac{r^2+r}{11}
\end{equation}
is strictly increasing for $r\ge2$.

Let $\rho=\rho_{11}$ and $\gamma=\rho-\beta$.  Then
\begin{equation}\label{eq:B11}
 B_{11}(r):=r^2M_{11}(r)
 =r^2\bigl(1+\cos(\alpha+\gamma)\bigr).
\end{equation}
Using $r=\sec\alpha$,
\[
 B_{11}(\alpha)
 =2\left(\frac{\cos((\alpha+\gamma)/2)}{\cos\alpha}\right)^2.
\]
Its logarithmic derivative has the sign of
\[
 2\tan\alpha-\tan\frac{\alpha+\gamma}{2}.
\]
Set
\[
 \Phi(\alpha)=\tan\frac{\alpha+\gamma}{2}-2\tan\alpha.
\]
On $\pi/3\le\alpha\le\alpha_*:=\pi-\rho$,
\[
 \frac{\alpha+\gamma}{2}\le\frac{5\pi}{12},
\]
and therefore
\[
 \Phi'(\alpha)
 =\frac12\sec^2\frac{\alpha+\gamma}{2}-2\sec^2\alpha
 \le(4+2\sqrt3)-8<0.
\]
At the right endpoint,
\[
 \frac{\alpha_*+\gamma}{2}=\frac{5\pi}{12}.
\]
Writing
\[
 c_*:=\cos\alpha_*=-x_{11}=\frac{20-11\sqrt3}{2}>0,
\]
the inequality $\Phi(\alpha_*)>0$ reduces, after squaring positive quantities,
to
\[
 (11+4\sqrt3)c_*^2-4
 =\frac{3097-1788\sqrt3}{4}>0.
\]
The last sign is exact because
\[
 3097^2-3\cdot1788^2=577>0.
\]
Hence $\Phi>0$ throughout the interval.  Since $\alpha=\arccos(1/r)$
is strictly increasing in $r$, $B_{11}$ is strictly decreasing on $[2,r_*]$.

Now define
\begin{equation}\label{eq:rstaragain}
 r_*:=\frac{40+22\sqrt3}{37}.
\end{equation}
Direct algebra gives
\begin{equation}\label{eq:reciprocal}
 \frac1{r_*}=10-\frac{11\sqrt3}{2}=-x_{11},
\end{equation}
so $\alpha_*+\rho=\pi$.  At this radius,
\[
 H_{11}(r_*)=1+\cos(\pi-\beta)=q,
\]
and
\[
 \frac{H_1(r_*)}{11}=\frac{1+1/r_*}{11}=q.
\]
Thus
\begin{equation}\label{eq:crossing}
 B_1(r_*)=B_{11}(r_*)=r_*^2q.
\end{equation}
For $2\le r<r_*$, the decreasing degree-eleven profile is strictly above its
value at $r_*$.  For $r>r_*$, the increasing degree-one profile is strictly above
its value at $r_*$.  This proves \cref{thm:optimality}.

The same exact algebra gives
\begin{equation}\label{eq:localidentity}
 r_*^2q\left(103-\frac{117\sqrt3}{2}\right)=1.
\end{equation}
Also $r_*>2$, since this is equivalent to $11\sqrt3>17$, whose square is
$363>289$.
